\section{Conclusão}
\label{sec:conclusion}

Os resultados mostram que nenhum banco de dados domina em todas as dimensões: o IoTDB lidera
em vazão de escrita sequencial e filtro de valor, o TimescaleDB oferece latência constante
em subamostragem e compatibilidade SQL completa, e o InfluxDB ocupa um meio-termo prático
entre desempenho e consumo de memória. As escolhas arquiteturais de cada sistema, e não o
ajuste de configuração, determinam o comportamento de escalonamento, e o isolamento de
recursos durante o benchmark afeta criticamente a validade dos resultados.

\subsection{Limitações}
\label{sec:limitations}

Os experimentos foram realizados em nó único, o que não captura a sobrecarga de rede real
de implantações distribuídas. Os testes de leitura compartilham o estado do cache de páginas
entre cargas de trabalho consecutivas, o que pode favorecer testes posteriores. A vazão de
escrita do IoTDB reflete reconhecimento em memória antes da persistência em disco;
configurações de durabilidade estrita podem reduzir esse valor. Por fim, o aquecimento da
JVM faz com que benchmarks curtos subestimem a vazão em estado estacionário do IoTDB.

\subsection{Trabalhos Futuros}
\label{sec:future-work}

Direções relevantes incluem avaliar implantações distribuídas (IoTDB multi-nó, InfluxDB
Clustered, TimescaleDB com Citus), testar com conjuntos de dados industriais de escala real
e reavaliar os compromissos após atualizações arquiteturais significativas, como o novo
motor colunar do InfluxDB~3.0. Medir o custo de replicação e failover também seria
essencial para decisões de implantação em produção.
